% ==================================================================
% id:       seq-019
% topic:    סדרות
% subtopic: הכרת סדרות: איברים ונוסחה לאיבר הכללי
% solution: full
% ==================================================================
האם
$42$
הוא איבר של הסדרה
$(a_n)_{n=1}^{\infty}$
כאשר
$a_n=n^2-5n$?
אם כן, מהו מיקומו בסדרה? ענו על אותה השאלה עם המספרים
$36$, $-4$
ו-$-6$.
\hint{לכל מספר פתרו את המשוואה שבה $a_n$ שווה לו, ובדקו אילו מהפתרונות הם מספרים טבעיים.}
\begin{solbox}
נזכיר כי
$$.a_n=n^2-5n$$
\begin{itemize}
\item
על מנת לענות על השאלה הראשונה, עלינו לוודא אם קיים פתרון (טבעי) למשוואה
$n^2-5n=42$.
נשים לב כי זוהי משוואה ריבועית וכי פתרונותיה של המשוואה הינם מספרים לא טבעיים, לכן התשובה לשאלה הראשונה היא שלילית.
\item
בדומה למקרה הקודם נפתור
$n^2-5n=36$.
לא קשה להשתכנע כי
$n=-4$
ו-$n=9$
הם הפתרונות שלה. לכן, התשובה לשאלה במקרה זה היא חיובית ומיקומו של
$36$
בסדרה הוא
$9$
(אין מיקום מינוס $4$).
\item
נמשיך עם
$-4$.
נפתור את המשוואה
$n^2-5n=-4$.
לא קשה להשתכנע כי
$n=1$
ו-$n=4$
הם הפתרונות של המשוואה. לכן
$-4$
הוא איבר של הסדרה והוא מופיע בה פעמיים, במקום הראשון ובמקום הרביעי.
\item
לבסוף, נפתור
$n^2-5n=-6$.
לא קשה להשתכנע כי
$n=2$
ו-$n=3$
הם הפתרונות שלה. לכן
$-6$
הוא איבר של סדרה שמופיע בה פעמיים: במקום השני ומיד אחריו.
\end{itemize}
\end{solbox}
